Topic 1: Semiconductor Physics (Intrinsic & Doping)
 Key Definitions
Intrinsic Semiconductor: A pure semiconductor crystal (like pure Silicon or Germanium) without any intentional impurities. Its electrical conductivity is solely determined by thermally generated electron-hole pairs.
Doping: The process of deliberately adding controlled amounts of impurity atoms (dopants) to an intrinsic semiconductor to significantly alter its electrical conductivity.
Pentavalent Impurity: An atom with 5 valence electrons (e.g., Phosphorus, Arsenic, Antimony). Used to create an n-type semiconductor, where electrons become the majority carriers.
Trivalent Impurity: An atom with 3 valence electrons (e.g., Boron, Aluminum, Indium). Used to create a p-typesemiconductor, where holes become the majority carriers.
 Essential Formulas & Variables
Formula Variables Breakdown
$n_i = N \cdot e^{-\frac{E_g}{2kT}}$ $n_i$: Intrinsic carrier concentration ($\text{m}^{-3}$)$N$: Semiconductor density factor / constant ($\text{m}^{-3}$)$E_g$: Bandgap energy gap ($\text{eV}$)$k$: Boltzmann’s constant ($8.62 \times 10^{-5}\text{ eV/K}$)$T$: Temperature in Kelvin ($K = ^\circ\text{C} + 273.15$)
$\sigma = q \cdot n_i \cdot (\mu_e + \mu_h)$ $\sigma$: Electrical conductivity ($\text{S/m}$ or $\Omega^{-1}\cdot\text{m}^{-1}$)$q$: Electron charge ($1.6 \times 10^{-19}\text{ C}$)$\mu_e$: Electron mobility ($\text{m}^2/\text{V}\cdot\text{s}$)$\mu_h$: Hole mobility ($\text{m}^2/\text{V}\cdot\text{s}$)
$n \cdot p = n_i^2$ (Mass Action Law) $n$: Free electron concentration$p$: Hole concentration
 Topic 2: Diode Characteristics & Resistance
 Key Definitions
Forward Bias: Connecting the positive terminal of a DC source to the anode (p-type) and the negative terminal to the cathode (n-type) of a diode to allow current flow.
Barrier Potential ($V_\gamma$): The internal voltage threshold required to overcome the depletion region. (Typically $0.7\text{ V}$ for Silicon, $0.3\text{ V}$ for Germanium).
Dynamic (AC) Resistance ($r_d$): The resistance offered by a diode to a changing/alternating signal current at a specific operating point. It is the reciprocal of the slope of the V-I curve.
 Essential Formulas & Variables
Formula Variables Breakdown
$r_d = \frac{\Delta V_d}{\Delta I_d}$ (Graphical) $r_d$: Dynamic / AC resistance ($\Omega$)$\Delta V_d$: Change in diode voltage between two points on a curve$\Delta I_d$: Change in diode current between those same two points
$r_d = \frac{26\text{ mV}}{I_d}$ (Theoretical) $26\text{ mV}$: Thermal voltage ($V_T$) at room temperature ($300\text{ K}$)$I_d$: DC operating forward current ($\text{mA}$) at the test point
 Topic 3: Power Supplies & Rectification (Filters)
 Key Definitions
Bridge Rectifier: A circuit configuration using four diodes that provides full-wave rectification, converting AC input into a pulsating DC output.
Ripple Voltage ($V_r$): The unwanted leftover AC variation riding on top of the rectified DC output voltage.
Peak Inverse Voltage (PIV): The maximum reverse-bias voltage that a diode can safely withstand when it is turned off without breaking down.
 Essential Formulas & Variables
Formula Variables Breakdown
$V_m = V_{\text{rms}} \cdot \sqrt{2}$ $V_m$: Peak input voltage ($\text{V}$)$V_{\text{rms}}$: Root-Mean-Square input voltage ($\text{V}$)
$V_r = \frac{V_m}{2 \cdot f \cdot R_L \cdot C}$(Approx)$V_r = \frac{V_m}{2f R_L C + 0.5}$ (Precise) $V_r$: Peak-to-peak ripple voltage ($\text{V}$)$f$: Line frequency ($\text{Hz}$) (Note: Ripple freq is $2f$ for full-wave)$R_L$: Connected load resistance ($\Omega$)$C$: Filter capacitance ($\text{F}$)
$V_{DC} = V_m - \frac{V_r}{2}$ $V_{DC}$: Final average DC load output voltage ($\text{V}$)
$\text{PIV} = V_m$ For full-wave bridge rectifiers, the required rating per diode.
